% !PRINTED 548
\weberpage{548}
%<<CONT
I. \gesperrt{The sum and difference of any two numbers in $\mathfrak{a}$ always yield numbers in $\mathfrak{a}$ again.}

II. \gesperrt{The product of any number in $\mathfrak{a}$ and a number in $\mathfrak{o}$ belongs to the system $\mathfrak{a}$.}

This requirement would be satisfied by the system consisting of the single number zero, which, however, for the sake of simplicity is not designated as an ideal.

The system $\mathfrak{o}$, on the other hand, is a proper ideal. Likewise the system $\mathfrak{o}\mu$ of all numbers in $\mathfrak{o}$ divisible by a definite number $\mu$ is an ideal, and such a one is called a \gesperrt{principal ideal}.

By the product $\mathfrak{a}\mathfrak{b}$ of two ideals $\mathfrak{a}$ and $\mathfrak{b}$ one understands the totality of all numbers which one obtains when one multiplies any number $\alpha$ from $\mathfrak{a}$ by a number $\beta$ from $\mathfrak{b}$ and adds an arbitrary number of such products of numbers, that is, the totality of all numbers of the form $\Sigma \alpha\beta$. That this product $\mathfrak{a}\mathfrak{b}$ is again an ideal is immediately evident. According to this definition, for example, $\mathfrak{o}\mathfrak{a} = \mathfrak{a}$, and the ideal $\mathfrak{o}$ plays the role of the unit in this multiplication.

One can now relate the ideals and the functionals to one another in such a way that the following laws prevail:

1. \gesperrt{To every integral functional there corresponds a definite ideal, and to associated functionals there corresponds the same ideal.}

2. \gesperrt{To every ideal there correspond infinitely many, but only associated, integral functionals.}

3. \gesperrt{To the product of two or more integral functionals there correspond the products of the ideals corresponding to the factors.}

4. \gesperrt{To an integer there corresponds a principal ideal.}

5. \gesperrt{To the units there corresponds the ideal $\mathfrak{o}$.}

To define this correspondence, we first assign to the system of all units the ideal $\mathfrak{o}$. If, further, $\varphi$ is any integral functional which is not a unit, then the totality of all integers $\alpha$ of the field $\Omega$ divisible by $\varphi$ satisfies, by \S.\,138, the requirements I, II, and is therefore an ideal, which we denote by $\mathfrak{a}$\footnote{For the notation of the ideals we use, with \gesperrt{Dedekind}, the small \gesperrt{German letters}.} and assign to the functional $\varphi$
%CONT>>
